\part{Transformations acido-basiques~; Étude d'une pile}
\begin{center}
    \textbf{Les deux parties sont indépendantes}
\end{center}
\section{Étude de l'ibuprofène comme acide carboxylique}
L'ibuprofène est une molécule de formule brute $\ce{C13H18O2}$. Elle constitue
le principe actif de divers médicaments de la classe des anti-inflammatoires.

Cette partie vise~:
\begin{itemize}
    \item l'étude d'une solution aqueuse d'ibuprofène~;
    \item le titrage d'une solution aqueuse d'ibuprofène.
\end{itemize}
\begin{donnees}
    $M(\ce{C13H18O2})=\qty{206}{\g\per\mol}$.
\end{donnees}

\subsection{Étude d'une solution aqueuse d'ibuprofène}
Le $\mathrm{pH}$ d'une solution aqueuse d'ibuprofène de concentration molaire
$C=\qty{5.0e-2}{\mol\per\L}$ vaut $\mathrm{pH}=2,7$ à
$\qty{25}{\degreeCelsius}$.

L'équation de la réaction modélisant la transformation entre l'ibuprofène et
l'eau est~:
\[
    \ce{C13H18O_{2(aq)} + H2O_{(\ell)} <=> C13H17O^-_{2(aq)} + H3O^+_{(aq)}}
\]

\begin{enumerate}
    \item Montrer que cette transformation est limitée.
    \item Calculer la valeur du quotient de réaction $Q_{r,\eq}$ du système
          chimique à l'équilibre.
    \item En déduire la valeur du $\mathrm{pK_{A}}$ du couple
          $(\ce{C13H18O_{2(aq)}}\ttslash{}\ce{C13H17O^-_{2(aq)}})$.
\end{enumerate}



\subsection{Titrage d'une solution aqueuse d'ibuprofène}
L'étiquette d'un médicament fournit l'information
<<\textbf{Ibuprofène.... 400 mg}>>.

On dissout un comprimé contenant l'ibuprofène selon un protocole bien défini
afin d'obtenir une solution aqueuse $\mathrm{(S)}$ d'ibuprofène de volume
$V_{S}=\qty{100}{\mL}$.

Pour vérifier, la masse d'ibuprofène contenu dans ce comprimé, on procède à un
titrage acido-basique du volume $V_{S}$ par une solution aqueuse d'hydroxyde de
sodium $\ce{Na^{+}_{(aq)} + HO^{-}_{(aq)}}$ de concentration molaire
$C_{B}=\qty{1.94e-1}{\mol\per\L}$, en utilisant le dispositif expérimental de
la Fig.~\ref{fig:2bac_svt_national_2018_normale_chim_1}.

La Fig.~\ref{fig:2bac_svt_national_2018_normale_chim_2} donne les courbes
$\mathrm{pH}=f(V_{B})$ et
$\frac{\mathrm{d}\mathrm{pH}}{\mathrm{d}V_{B}}=g(V_{B})$ obtenues lors de ce
dosage.

\begin{minipage}{0.49\linewidth}
    \begin{figure}[H]
        \centering
        \includegraphics[width=0.75\textwidth]{2026-03-12_015233-}
        \caption{}
        \label{fig:2bac_svt_national_2018_normale_chim_1}
    \end{figure}
\end{minipage}
\hfill
\begin{minipage}{0.49\linewidth}
    \begin{figure}[H]
        \centering
        \includegraphics[width=0.8\textwidth]{2026-03-12_015243-}
        \caption{}
        \label{fig:2bac_svt_national_2018_normale_chim_2}
    \end{figure}
\end{minipage}

\begin{enumerate}
    \item Nommer les éléments du dispositif expérimental numérotés 1, 2, 3 et 4
          sur la Fig.~\ref{fig:2bac_svt_national_2018_normale_chim_1}.
    \item Parmi les courbes~(1) et~(2) de la
          Fig.~\ref{fig:2bac_svt_national_2018_normale_chim_2}, quelle est celle
          qui représente $\mathrm{pH}=f(V_{B})$~?
    \item Déterminer graphiquement la valeur du volume $V_{B,E}$ versé à
          l'équivalence.
    \item Écrire l'équation de la réaction qui a eu lieu lors du dosage sachant
          qu'elle est totale.
    \item Calculer la valeur de la quantité de matière $n_{A}$ d'ibuprofène dans
          la solution $\mathrm{(S)}$.
    \item Déduire la valeur de la masse $m$ d'ibuprofène dans le comprimé et la
          comparer à celle indiquée sur l'étiquette du médicament.
\end{enumerate}



\section{Étude d'une pile}
Les piles constituent des systèmes chimiques dont le fonctionnement est basé sur
des réactions d'oxydo-réductions. L'étude de ces systèmes permet de prévoir le
sens de leur évolution et reconnaitre le fonctionnement de ces piles.

Cette partie vise la détermination de la durée de fonctionnement de la pile
(Zinc/Cuivre) schématisée dans la figure ci-contre.

\begin{minipage}{0.69\linewidth}
    \begin{donnees}
        \begin{itemize}
          \item Masse de la partie immergée de l'électrode de Zinc~:
                $m=\qty{6,54}{\gram}$~;
          \item Volume de chaque solution~: $V=\qty{50}{\mL}$~;
          \item Concentration de chaque solution~: $C=\qty{1,0}{\M}$~;
          \item $1~\mathcal{F}=\qty{9,65e4}{\coulomb\per\mole}$~;
          \item $M(\ce{Zn})=\qty{65,4}{\gram\per\mole}$.
        \end{itemize}
    \end{donnees}
    On laisse fonctionner la pile pendant une durée $\Delta t$ suffisamment longue
    jusqu'à ce que la pile ne débite plus.
    
    L'équation bilan du fonctionnement de cette pile est~:
    \[
        \ce{Zn_{(s)} + Cu^{2+}_{(aq)} -> Zn^{2+}_{(aq)} + Cu_{(s)}}
    \]
\end{minipage}
\hfill
\begin{minipage}{0.30\linewidth}
    \begin{figure}[H]
        \centering
        \begin{tikzpicture}
            \newdimen\bh \bh=1.2cm
            \foreach \x/\c [count=\i] in {0/capri,3.5/lavendergray}{%
            \node[inner sep=0pt,minimum width=2cm,
                minimum height=\bh] (b\i) at (\x,0) {};
            \fill[\c,opacity=0.5,rounded corners]
                (b\i.north west) |- (b\i.south east) [sharp corners] --
                (b\i.north east) -- cycle;
            \draw[very thick,rounded corners,shorten <=-5mm,shorten >=-5mm]
                (b\i.north west) |- (b\i.south east) -- (b\i.north east);
            }
            % pont salin
            \coordinate (pb) at
                ([xshift=-5mm]$(b1.south east)!.4!(b1.north east)$);
            \coordinate (pe) at
                ([xshift=5mm]$(b2.south west)!.4!(b2.north west)$);
            \draw[double,thick,double distance=2pt,rounded corners]
                (pb) -- ++(0,1.5) coordinate(tmp) -- node[inner sep=2pt] (ps) {}
                (tmp -| pe) --(pe);
            \node[above,align=center] (pss) at (ps) {pont salin};
            % électrodes
            \node[inner sep=0pt,minimum width=4mm,minimum height=2cm,
                draw,thick,left color=white,right color=copper,
                anchor=south] (e1) at ([xshift=-1cm]pb) {};
            \node[inner sep=0pt,minimum width=4mm,minimum height=2cm,
                draw,thick,left color=white,right color=gray,
                anchor=south] (e2) at ([xshift=1cm]pe) {};
            % Text
            \draw[{Circle[scale=0.7]}-,shorten <=-5mm] (b1.south) -- ++(0,-0.3)
                node[below] {$\ce{Cu^{2+}_{(aq)} + SO$^{2-}_{4(aq)}$}$};
            \draw[{Circle[scale=0.7]}-,shorten <=-5mm] (b2.south) -- ++(0,-0.3)
                node[below] {$\ce{Zn^{2+}_{(aq)} + SO$^{2-}_{4(aq)}$}$};
            \node[below left,align=center] at (e1.north west) {$\ce{Cu}$};
            \node[below right,align=center] at (e2.north east) {$\ce{Zn}$};
            % Pôles
            %\node[circle,inner sep=4pt,draw,thick] (ca) at
            %    ([shift=(145:0.5)]e1.north west) {};
            %\draw[thick] (ca.east) -- (ca.west) (ca.north) -- (ca.south);
            %\node[circle,inner sep=4pt,draw,thick] (an) at
            %    ([shift=(35:0.5)]e2.north east) {};
            %\draw[thick] (an.east) -- (an.west);
            % Fils
            \draw (e1.north) to[short,i=$I$] ++(0,1.5)
                to[R,l=$R$] ++(2.4,0) coordinate(tmp)
                to[rmeter,t=A] (tmp -| e2.north) to[short] (e2.north);
        \end{tikzpicture}
    \end{figure}
\end{minipage}

\begin{enumerate}
    \item Recopier sur votre copie le numéro de la question et écrire la lettre
          correspondante à la proposition vraie.

          Le schéma conventionnel de cette pile est~:
          \begin{enumerate}[label={\textbf{\Alph*)}}]
              \item $\circletlineh~\ce{Cu_{(s)}}~|~\ce{Cu^2+_{(aq)}}~\parallel~
                    \ce{Zn^2+_{(aq)}}~|~\ce{Zn_{(s)}}~\circletlinevh$
              \item $\circletlineh~\ce{Zn_{(s)}}~|~\ce{Zn^2+_{(aq)}}~\parallel~
                    \ce{Cu^2+_{(aq)}}~|~\ce{Cu_{(s)}}~\circletlinevh$
              \item $\circletlinevh~\ce{Zn_{(s)}}~|~\ce{Zn^2+_{(aq)}}~\parallel~
                    \ce{Cu^2+_{(aq)}}~|~\ce{Cu_{(s)}}~\circletlineh$
              \item $\circletlinevh~\ce{Cu^2+_{(aq)}}~|~\ce{Cu_{(s)}}~\parallel~
                    \ce{Zn_{(s)}}~|~\ce{Zn^2+_{(aq)}}~\circletlineh$
          \end{enumerate}
    \item Montrer que la quantité de matière du cuivre déposé est
          $n(\ce{Cu})=\qty{5e-2}{\mole}$.
    \item Déterminer la valeur de la durée $\Delta t$ du fonctionnement de la
          pile sachant qu'elle délivre un courant d'intensité constante
          $I=\qty{100}{\mA}$.
\end{enumerate}


